\section{Lecture 2: Light Condensed Sets and Abelian Groups}
\label{lec:2}

This lecture, the first given by Scholze, develops the foundation on which the
whole course rests: \emph{light condensed sets} and, on top of them, light
condensed abelian groups. After some motivating remarks on what one might
ultimately want analytic geometry to do, Scholze recalls the category of
profinite sets and Stone duality, isolates the \emph{light} profinite sets by a
countability condition, establishes the two special properties of light
profinite sets that make the theory run, defines light condensed sets and
compares the construction with topological spaces (and with earlier
topos-theoretic proposals of Johnstone and of Escard\'o--Xu), and finally passes
to light condensed abelian groups, whose good homological algebra is the point
of the whole exercise.

\subsection{Motivation: a global perfectoid space}

Scholze framed the course by a personal goal of some ten years' standing. Much
recent $p$-adic geometry --- perfectoid spaces, prismatic cohomology, the
Fargues--Fontaine curve and the geometrization of the local Langlands
correspondence \cite{fargues2024geometrizationlocallanglandscorrespondence} ---
works beautifully with the $p$-adic numbers, but at the cost of quite fancy
$p$-adic geometry. One would hope that similar techniques apply not only to
$p$-adic local fields but also to real local fields --- something Scholze thinks
is now genuinely possible --- and then also to the integers, globally. What one
would like is some notion of a \emph{global} object: an analytic space, modelled
by rings carrying a nontrivial topology, incorporating both non-Archimedean and
Archimedean parts, and extremely non-Noetherian, so that none of the usual
finiteness properties can be imposed.

\begin{remark}[A vague, possibly misguided idea]
\label{rmk:2:global-perfectoid}
Scholze stressed repeatedly that the following is only a vague idea, quite
possibly completely misguided. In the perfectoid world one adjoins to a perfect
coefficient field a variable together with all of its $p$-power roots; the
prototypical example is
\begin{equation}
\label{eq:2:perfectoid-prototype}
\mathbb{C}_p\bigl\langle T^{\pm 1/p^{\infty}}\bigr\rangle
=\mathbb{C}_p\bigl\langle T^{\mathbb{Z}[1/p]}\bigr\rangle .
\end{equation}
To imitate this globally one wants to get rid of the fixed prime, so one is led
to adjoin a variable with \emph{all} of its rational powers, and --- as the
analogue at the infinite place --- with all of its real powers:
\begin{equation}
\label{eq:2:global-idea}
K\bigl\langle T^{\mathbb{Q}}\bigr\rangle
\subset
K\bigl\langle T^{\mathbb{R}}\bigr\rangle .
\end{equation}
It is unclear what kind of object $K\langle T^{\mathbb{R}}\rangle$ should be. One
can of course treat $\mathbb{R}$ as discrete and simply adjoin the real powers,
getting some algebra; but this feels artificial (it presents $\mathbb{R}$ as an
uncountable-dimensional $\mathbb{Q}$-vector space and does not improve the
situation). One really wants to remember the topology on $\mathbb{R}$; but then
one is mixing the $p$-adic-type topology, where these objects sit on a Banach
line, with the real topology --- so this is definitely \emph{not} a Banach
algebra. It is unclear technically what it should be, and still less clear what
geometric object should correspond to it. These may not be the right objects at
all; the point is to have a language in which one can at least talk about them.
\end{remark}

The reaction to these difficulties is the condensed formalism. When Clausen came
to Bonn in 2018 he suggested replacing topological spaces by condensed sets, and
since then the program has been to replace
\[
\text{topological spaces / groups / rings}
\quad\rightsquigarrow\quad
\text{condensed sets / groups / rings.}
\]
This switch resolves many of the foundational problems of analytic geometry ---
for instance it removes the Noetherian hypotheses pervasive in the theory of
adic spaces and in the associated derived categories of complete modules --- and
it makes objects such as $K\langle T^{\mathbb{R}}\rangle$ at least \emph{exist},
now as condensed rings. The attitude of the course is therefore to develop the
most natural possible foundations for analytic geometry based on condensed rings,
and to try them out on the known formalisms (allowing, along the way, such
variants as overconvergent function theories), in the hope of eventually
accommodating more exotic examples of the kind above.

\subsection{Profinite sets and Stone duality}

\begin{notation}
\label{not:2:light-drop}
Throughout, ``ring'' means commutative ring, and (following Lecture~1) the word
\emph{light} is eventually dropped: from the point where light condensed sets are
defined, ``condensed set'' means light condensed set and ``profinite set'' means
light profinite set.
\end{notation}

The starting point is the category $\ProFin$ of profinite sets, which can be
described in three equivalent ways.

\begin{proposition}[Stone duality]
\label{prop:2:stone-duality}
The following categories are equivalent:
\begin{enumerate}
\item $\mathrm{Pro}(\mathrm{Fin})$, the pro-category of finite sets. Its objects
are formal cofiltered inverse limits ``$\varprojlim_{i\in I} S_i$'' with each
$S_i$ finite and $I$ a cofiltered poset (nonempty, and for all $i,j\in I$ there is
$k$ with $k\le i,j$); its morphisms are
\begin{equation}
\label{eq:2:pro-hom}
\Homop\Bigl(\text{``}\varprojlim_{i\in I} S_i\text{''},\
\text{``}\varprojlim_{j\in J} T_j\text{''}\Bigr)
=\varprojlim_{j}\ \varinjlim_{i}\ \Homop(S_i,T_j).
\end{equation}
\item Totally disconnected compact Hausdorff spaces.
\item The opposite of the category of Boolean algebras, a Boolean algebra being a
commutative ring $R$ with $x^2=x$ for all $x\in R$.
\end{enumerate}
The equivalences go
\[
\text{``}\varprojlim_i S_i\text{''}\ \longmapsto\
S=\varprojlim_i S_i\ \longmapsto\
\Cont(S,\mathbb{F}_2)=\varinjlim_i \mathbb{F}_2^{S_i},
\]
where $S$ is the actual inverse limit (a compact Hausdorff space) and
$\Cont(S,\mathbb{F}_2)$ is its Boolean algebra of continuous
$\mathbb{F}_2$-valued functions; conversely a Boolean algebra $A$ is sent to
$\Homop(A,\mathbb{F}_2)=\Specop A$, its set of points, which is canonically the
inverse limit $\varprojlim_{A_j\subseteq A} \Homop(A_j,\mathbb{F}_2)$ over the
finite subalgebras $A_j$ (as $A=\varinjlim_j A_j$ is their filtered union, so
$\Homop(A,\mathbb{F}_2)=\varprojlim_j\Homop(A_j,\mathbb{F}_2)$ exhibits $\Specop A$
as a profinite set).
\end{proposition}

Note that in a Boolean algebra $(-1)^2=-1$, whence $2=0$: these are
$\mathbb{F}_2$-algebras. Scholze remarked that he will usually think of a
profinite set through the presentation of \Cref{prop:2:stone-duality}(1), as a
formal inverse limit of finite sets.

\subsubsection*{Two ways to measure the size of a profinite set}

\begin{definition}[Size and weight]
\label{def:2:size-weight}
Let $S=\varprojlim_i S_i$ be a profinite set.
\begin{enumerate}
\item The \emph{size} of $S$ is the cardinality of its underlying set,
$\kappa:=|S|$.
\item The \emph{weight} of $S$ is $\lambda:=\bigl|\Cont(S,\mathbb{F}_2)\bigr|$,
the cardinality of the associated Boolean algebra.
\item $S$ is \emph{light} if $\lambda\le\omega$, i.e.\ if its weight is at most
countable.
\end{enumerate}
\end{definition}

\begin{remark}
\label{rmk:2:weight-index}
If $\lambda$ is infinite it equals $|I|$ for the smallest possible index poset
$I$ in a presentation $S=\varprojlim_{i\in I}S_i$ (a point, say, can be presented
with an arbitrarily large index poset, but there is always a minimal choice). In
particular a profinite set is light if and only if it can be written as a
\emph{countable} (sequential) inverse limit of finite sets.
\end{remark}

\begin{example}[Size vs.\ weight]
\label{ex:2:size-weight}
\begin{enumerate}
\item The one-point compactification $S=\mathbb{N}\cup\{\infty\}=\varprojlim_n
\{0,1,\dots,n,\infty\}$ (transition maps collapsing everything from $n$ onward to
$\infty$): $\kappa=\omega$, $\lambda=\omega$. This is the smallest infinite
profinite set.
\item The Cantor set $S=\{0,1\}^{\mathbb{N}}=\varprojlim_n\{0,1\}^n$:
$\kappa=2^{\omega}$, $\lambda=\omega$. It is much bigger than
$\mathbb{N}\cup\{\infty\}$ in cardinality of points, but still light.
\item The Stone--\v{C}ech compactification $S=\beta\mathbb{N}$: here
\[
\Cont(\beta\mathbb{N},\mathbb{F}_2)=\Cont(\mathbb{N},\mathbb{F}_2)
=\{\text{subsets of }\mathbb{N}\},
\]
by the universal property of $\beta\mathbb{N}$ and discreteness of $\mathbb{F}_2$,
so $\lambda=2^{\omega}$ and $\kappa=2^{2^{\omega}}$. This is not light.
\end{enumerate}
Locally compact spaces have two compactifications: the one-point compactification,
which is small, and the Stone--\v{C}ech compactification, which is huge.
\end{example}

\begin{proposition}[Comparison of size and weight]
\label{prop:2:size-weight-ineq}
For any profinite set,
\begin{equation}
\label{eq:2:size-weight-ineq}
\lambda\le 2^{\kappa},\qquad \kappa\le 2^{\lambda}.
\end{equation}
Moreover, if $\kappa=\omega$ then $\lambda=\omega$.
\end{proposition}

\begin{proof}
For the first inequality, $\lambda=|\Cont(S,\mathbb{F}_2)|\le
|\mathrm{Map}(S,\mathbb{F}_2)|=2^{\kappa}$. For the second, writing $A$ for the
Boolean algebra (so $|A|=\lambda$),
$\kappa=|\Homop(A,\mathbb{F}_2)|\le|\mathrm{Map}(A,\mathbb{F}_2)|=2^{\lambda}$.
For the last statement, enumerate $S=\{s_0,s_1,s_2,\dots\}$. Inductively choose,
for each $n$, a finite quotient $S\twoheadrightarrow S_n$ (compatibly with
$S_n\to S_{n-1}$) into which $\{s_0,\dots,s_n\}$ injects. Then $S\to
\varprojlim_n S_n$ is injective (any two points are already separated in some
finite quotient) and surjective (being a sequential limit of surjections), hence
a bijection, exhibiting $S=\varprojlim_n S_n$ with countable weight.
\end{proof}

\begin{remark}[Weight versus size]
\label{rmk:2:gabber}
Scholze noted that the same proof in fact shows more: if $\kappa$ is infinite
then $\lambda\le\kappa$. Indeed, since $\kappa^2=\kappa$, one can choose for each
of the $\kappa$ pairs of distinct points of $S$ a finite quotient separating
them; these generate a point-separating subalgebra of $\Cont(S,\mathbb{F}_2)$ of
size at most $\kappa$, which must then be all of it. Thus in the infinite case the
weight never exceeds the size. The bound $\kappa\le 2^{\lambda}$, on the other
hand, is attained (e.g.\ for the Cantor set and $\beta\mathbb{N}$, where
$\kappa=2^{\lambda}$). Scholze had at first suspected that $\lambda>\kappa$ might
be achievable by some complicated example, but this improved bound rules that out
--- whatever that example computed, it was something else. For us only the
elementary facts of \Cref{prop:2:size-weight-ineq} are needed.
\end{remark}

Specializing Stone duality to the light case gives:

\begin{proposition}[Stone duality, light version]
\label{prop:2:stone-duality-light}
The following categories are equivalent:
\begin{enumerate}
\item $\mathrm{Pro}_{\mathbb{N}}(\mathrm{Fin})$, the sequential pro-category of
finite sets, whose objects are formal inverse limits
``$\varprojlim_{n\in\mathbb{N}} S_n$'' along the integers with each $S_n$ finite,
and whose morphisms are
\begin{equation}
\label{eq:2:pro-N-hom}
\Homop\Bigl(\text{``}\varprojlim_n S_n\text{''},\
\text{``}\varprojlim_m T_m\text{''}\Bigr)
=\varprojlim_m\ \varinjlim_n\ \Homop(S_n,T_m)
=\colimop_{f:\mathbb{N}\to\mathbb{N}\ \mathrm{str.\ incr.}}\ \varprojlim_m
\Homop\bigl(S_{f(m)},T_m\bigr).
\end{equation}
\item Metrizable totally disconnected compact Hausdorff spaces.
\item The opposite of the category of countable Boolean algebras.
\end{enumerate}
\end{proposition}

The second description of morphisms in \eqref{eq:2:pro-N-hom} uses that in a
pro-object one may always pass to a cofinal subsequence: a morphism is a strictly
increasing reindexing $f$ of the source together with a compatible family of maps
$S_{f(m)}\to T_m$. The metrizability in (2) is exactly the reformulation, for
compact Hausdorff spaces, of the sequential/compatibility condition in (1).

\begin{proposition}[Countable limits]
\label{prop:2:countable-limits}
The category of light profinite sets has all countable limits. Moreover a
sequential limit of surjections is surjective.
\end{proposition}

Here surjectivity is meant on underlying point sets (equivalently on the
underlying compact Hausdorff spaces). In the sequential case this is elementary
--- one successively lifts --- whereas the corresponding statement for all
profinite sets requires a rather high-powered compactness argument.

\subsection{Two properties special to light profinite sets}

\begin{proposition}[Cantor surjections are universal]
\label{prop:2:cantor-surjection}
Let $S$ be a light profinite set. Then there exists a surjection from the Cantor
set,
\begin{equation}
\label{eq:2:cantor-surjection}
\{0,1\}^{\mathbb{N}}\twoheadrightarrow S.
\end{equation}
\end{proposition}

\begin{proof}
A simple induction: write $S=\varprojlim_n S_n$ and at each stage pick a large
enough finite power $\{0,1\}^{k_n}$ surjecting onto $S_n$ compatibly. (Exercise.)
\end{proof}

The next two properties genuinely single out the light profinite sets; they may
hold for some other profinite sets, but that would be hard to characterize, and
they are among the key technical reasons for making the switch to the light
setting.

\begin{proposition}[Open subsets are tame]
\label{prop:2:open-subsets}
Let $S$ be a light profinite set and $U\subseteq S$ open. Then $U$ is a
\emph{countable disjoint union} of light profinite sets.
\end{proposition}

\begin{proof}
Write $S=\varprojlim_n S_n$, let $Z=S\setminus U$ be the closed complement,
$Z=\varprojlim_n Z_n$ with $Z_n=\mathrm{im}(Z\to S_n)\subseteq S_n$, and
$\pi_n\colon S\to S_n$ the projections. Then
\[
U=\bigcup_n \pi_n^{-1}(S_n\setminus Z_n),
\]
each $\pi_n^{-1}(S_n\setminus Z_n)\subseteq S$ being clopen (hence itself light
profinite), so $U$ is an increasing union of clopens. Writing it as the
associated disjoint union,
\begin{equation}
\label{eq:2:open-disjoint-union}
U=\coprod_n\Bigl(\pi_{n+1}^{-1}(S_{n+1}\setminus Z_{n+1})\setminus
\pi_n^{-1}(S_n\setminus Z_n)\Bigr),
\end{equation}
a countable disjoint union of light profinite sets.
\end{proof}

\begin{remark}[Wildness in the general case]
\label{rmk:2:general-open-wild}
This is to be contrasted with the general (non-light) case: there exists a
profinite set $S$ with an open subset $U\subseteq S$ that is \emph{not} a disjoint
union of profinite sets, and indeed with $H^i(U,\mathbb{Z})\neq 0$ for some
$i>0$. Disjoint unions take cohomology to products, and profinite sets are
totally split (no nontrivial covers, so higher cohomology vanishes), so such
higher cohomology on $U$ witnesses the failure. In the light case, by contrast,
\Cref{prop:2:open-subsets} gives complete control. In discussion an audience
member gave concrete non-light counterexamples to both this proposition and
\Cref{prop:2:injective}, built from the ordinal $[0,\omega_1]$ (a profinite set):
its complement of the top point $\omega_1$ is an open subset that is not a
disjoint union of clopens, and a suitable closed subset of $[0,\omega_1]^2$ is
not a retract of the product. Scholze added that one can also produce opens with
nonvanishing higher cohomology, e.g.\ from $\beta\mathbb{N}$ by removing boundary
points; he thought the first infinite profinite set with a removed limit point
also works but was not certain in that case.
\end{remark}

\begin{proposition}[Injectivity]
\label{prop:2:injective}
Let $S$ be a light profinite set. Then $S$ is an injective object in
$\mathrm{Pro}(\mathrm{Fin})$: for every injection $Z\hookrightarrow X$ of
profinite sets, every map $Z\to S$ extends to $X$.
\end{proposition}

That is, one can always fill in the dashed arrow in
% board 17 @ 00:56:23 — /home/deven/documents/coding/vms/gpu/home/notetaker/output/peter-scholze-224-analytic-stacks/boards/board-17.jpg
\[
\begin{tikzcd}
Z \arrow[r] \arrow[d, hook] & S \\
X \arrow[ur, dashed, "\exists"']
\end{tikzcd}
\]

\begin{proof}
First take $S=\{0,1\}=\mathbb{F}_2$. Then the claim is that
$\Cont(X,\mathbb{F}_2)\to\Cont(Z,\mathbb{F}_2)$ is surjective, i.e.\ any clopen
subset of $Z$ extends to a clopen subset of $X$. In general write
$S=\varprojlim_n S_n$ with all $S_{n+1}\to S_n$ surjective and induct on $n$:
having extended the map to $S_n$, decompose the situation into the disjoint union
over the fibres of $S_{n+1}\to S_n$, reducing to the case $S_n=\ast$ and hence to
$S_{n+1}$ finite, which follows from the $\mathbb{F}_2$-case.
\end{proof}

\begin{remark}
\label{rmk:2:injective-more}
These light profinite sets do not exhaust the injective objects of
$\mathrm{Pro}(\mathrm{Fin})$: injective objects are closed under products, so any
product of light profinite sets is again injective, and there are much larger
ones. It is an instructive exercise to see exactly where the inductive argument
above breaks down for a general (non-sequential) profinite set. Scholze also
raised the edge case: if $Z$ is empty but $X$ is not, one must know $S\neq
\emptyset$ for the extension to exist.
\end{remark}

\subsection{Light condensed sets}

\begin{definition}[Light condensed set]
\label{def:2:light-condensed-set}
A \emph{light condensed set} is a sheaf on the category of light profinite sets
for the Grothendieck topology generated by
\begin{enumerate}
\item finite disjoint unions, and
\item all surjective maps.
\end{enumerate}
Equivalently, it is a functor
$X\colon\mathrm{Pro}_{\mathbb{N}}(\mathrm{Fin})^{\mathrm{op}}\to\mathrm{Sets}$,
$S\mapsto X(S)$, such that
\begin{enumerate}
\item[(0)] $X(\emptyset)=\ast$;
\item[(1)] $X(S_1\amalg S_2)\xrightarrow{\ \sim\ }X(S_1)\times X(S_2)$;
\item[(2)] for every surjection $T\xrightarrow{f}S$,
\begin{equation}
\label{eq:2:sheaf-equalizer}
X(S)\xrightarrow{\ \sim\ }\eq\Bigl(X(T)
\underset{p_2^{*}}{\overset{p_1^{*}}{\rightrightarrows}}
X(T\times_S T)\Bigr),
\end{equation}
where $p_1,p_2\colon T\times_S T\to T$ are the two projections.
\end{enumerate}
\end{definition}

Condition (2) is exactly what the sheaf condition for the surjective cover $f$
unwinds to: a map $S\to X$ is determined by its restriction to $T$, and the maps
$T\to X$ that descend to $S$ are those whose two pullbacks to $T\times_S T$ agree
(they ``agree on fibres'' of $f$).

\begin{example}[Topological spaces]
\label{ex:2:top-to-condensed}
Any topological space $A$ gives a light condensed set
\begin{equation}
\label{eq:2:underline-A}
\underline{A}\colon S\longmapsto \Cont(S,A),
\end{equation}
the continuous maps from the compact Hausdorff space $S$ into $A$, with all its
functoriality. That condition (2) holds is not a complete tautology: it says that
a set-theoretic map $S\to A$ which becomes continuous after restriction along a
surjection $T\twoheadrightarrow S$ was already continuous --- equivalently, that
$S$ carries the quotient topology from $T$, which holds because surjections of
compact Hausdorff spaces are quotient maps. Evaluating,
\[
\underline{A}(\ast)=A\ \text{as a set},\qquad
\underline{A}(\mathbb{N}\cup\{\infty\})=\{\text{convergent sequences in }A,\
\text{with a choice of limit point}\},
\]
and $\underline{A}(\{0,1\}^{\mathbb{N}})$ (functions on the Cantor set) carries an
action of $\Endop(\{0,1\}^{\mathbb{N}})$.
\end{example}

For a general light condensed set $X$ one thinks of $X(\ast)$ as the underlying
set and of $X(\mathbb{N}\cup\{\infty\})$ as the ``convergent sequences in $X$'',
where giving such a convergent sequence includes giving a witness for its limit
(and there may be several limit points).

\begin{remark}[A light condensed set is almost a monoid action]
\label{rmk:2:monoid-action}
A light condensed set $X$ is completely determined by the single set
$X(\{0,1\}^{\mathbb{N}})$ together with the action on it of the abstract monoid
$\Endop(\{0,1\}^{\mathbb{N}})$:
\[
\underbrace{X(\{0,1\}^{\mathbb{N}})}_{\text{abstract set}}\ \circlearrowleft\
\underbrace{\Endop(\{0,1\}^{\mathbb{N}})}_{\text{abstract monoid}}.
\]
Indeed, every profinite set is covered by a Cantor set
(\Cref{prop:2:cantor-surjection}), and the relevant fibre products are again so
covered, so all values $X(S)$ are recovered from $X(\{0,1\}^{\mathbb{N}})$ via the
sheaf condition, the transition maps being induced by endomorphisms of the Cantor
set. One could therefore describe light condensed sets as certain sets with an
action of this (very large) abstract monoid, but Scholze regards this as a
strictly worse way to think about them: the functor from light condensed sets to
$\Endop(\{0,1\}^{\mathbb{N}})$-sets is fully faithful but not essentially
surjective --- the covering (sheaf) condition cuts out which actions arise, and
that condition is not visible from the monoid alone. The point of the remark is
that a light condensed set is nevertheless a very concrete, almost combinatorial
object.
\end{remark}

\subsubsection*{Topological realization}

\begin{proposition}[Left adjoint to $A\mapsto\underline{A}$]
\label{prop:2:top-realization}
The functor $\mathrm{Top}\to\CondSetl$, $A\mapsto\underline{A}$, has a left
adjoint $X\mapsto X(\ast)_{\mathrm{top}}$, where $X(\ast)_{\mathrm{top}}$ is the
underlying set $X(\ast)$ equipped with the quotient topology from
\begin{equation}
\label{eq:2:top-realization}
\coprod_{\substack{S\ \mathrm{light\ profinite}\\ \alpha\in X(S)}} S\
\longrightarrow\ X(\ast),
\end{equation}
the map sending a point of the copy of $S$ indexed by $\alpha\in X(S)$ to the
corresponding point of $X(\ast)$. The resulting space $X(\ast)_{\mathrm{top}}$ is
metrizably compactly generated.
\end{proposition}

Here a topological space is \emph{metrizably compactly generated} if continuity
of a map out of it may be tested on maps into it from metrizable compact
Hausdorff spaces; by construction of the quotient topology one only has to test
continuity on the (metrizable) Cantor sets appearing in
\eqref{eq:2:top-realization}.

\begin{corollary}[Topological spaces as condensed sets]
\label{cor:2:top-fully-faithful}
The functor $A\mapsto\underline{A}$ restricts to a fully faithful embedding
\[
\{\text{metrizably compactly generated topological spaces}\}\hookrightarrow
\CondSetl,
\]
and for any metrizably compactly generated space $A$ the counit is an
isomorphism, $\underline{A}(\ast)_{\mathrm{top}}\xrightarrow{\ \sim\ }A$.
\end{corollary}

This is a very weak condition: virtually all topological spaces arising in
practice --- in particular all the topological modules, Banach spaces and the
like used in condensed functional analysis --- are metrizably compactly
generated. The condition is even implied by being metrizable, or more generally
sequential.

\begin{remark}[Earlier topos-theoretic proposals]
\label{rmk:2:johnstone-escardo}
Two related constructions have appeared before.
\begin{enumerate}
\item Johnstone's ``topological topos'' \cite{Johnstone_1979} is built on just the
single test object $\mathbb{N}\cup\{\infty\}$ (the sequence space), with the
\emph{canonical} topology (the finest for which all representable presheaves are
sheaves). This topology is not finitary: it admits genuinely infinite covers with
no finite subcover. Much of what Johnstone does goes through, but the
non-finitary topology leads to bad algebraic behaviour.
\item Escard\'o--Xu \cite{Escard__2016} use light profinite sets but with a
finitary topology, taking only finite disjoint unions as covers --- \emph{not}
all surjective maps. Allowing all surjective maps as covers is, however,
important for the good algebraic properties of the theory (as will be seen for
abelian groups below), and this is precisely the point at which the present
setup differs from theirs.
\end{enumerate}
\end{remark}

\subsection{Light condensed abelian groups}

Piling abelian-group structure on top gives the category
$\CondAbl$ of light condensed abelian groups: sheaves of abelian groups on the
site of light profinite sets.

\begin{proposition}[Grothendieck abelian category]
\label{prop:2:grothendieck}
For sheaves on any site, sheaves of abelian groups form a Grothendieck abelian
category: all limits and colimits exist and filtered colimits are exact.
In particular $\CondAbl$ is a Grothendieck abelian category.
\end{proposition}

The whole point of passing to condensed abelian groups, as recalled in Lecture~1,
is that dense inclusions --- which make categories of topological abelian groups
fail to be abelian --- are now handled correctly. Scholze illustrated this with
the two maps
\begin{equation}
\label{eq:2:dense-inclusions}
\underline{\mathbb{Q}}\hookrightarrow\underline{\mathbb{R}},
\qquad
\underline{\Rdisc}\to\underline{\mathbb{R}},
\end{equation}
where $\mathbb{Q}$ carries its usual topology and $\Rdisc$ is $\mathbb{R}$ with
the discrete topology. Both are perfectly nice maps of topological abelian
groups whose cokernels are problematic topologically; condensed mathematics
computes them cleanly.

\begin{example}[Two cokernels]
\label{ex:2:cokernels}
For the first map, the cokernel $\underline{\mathbb{R}}/\underline{\mathbb{Q}}$
has
\begin{equation}
\label{eq:2:R-mod-Q}
\bigl(\underline{\mathbb{R}}/\underline{\mathbb{Q}}\bigr)(\ast)=\mathbb{R}/\mathbb{Q},
\qquad
\bigl(\underline{\mathbb{R}}/\underline{\mathbb{Q}}\bigr)(S)
=\Cont(S,\mathbb{R})\big/\Cont(S,\mathbb{Q}),
\end{equation}
and one can prove that this naive quotient presheaf is already a sheaf, so no
sheafification is needed: the value at $\ast$ is the plain quotient, while the
values at general $S$ remember the topology (continuous $\mathbb{R}$-valued
functions modulo continuous $\mathbb{Q}$-valued ones).

For the second map the phenomenon is more dramatic. The underlying set is
\begin{equation}
\label{eq:2:R-mod-Rdisc-point}
\bigl(\underline{\mathbb{R}}/\underline{\Rdisc}\bigr)(\ast)=\mathbb{R}/\mathbb{R}=0,
\end{equation}
yet the cokernel is not zero: on a general profinite set
\begin{equation}
\label{eq:2:R-mod-Rdisc-general}
\bigl(\underline{\mathbb{R}}/\underline{\Rdisc}\bigr)(S)
=\Cont(S,\mathbb{R})\big/\Cont(S,\Rdisc),
\end{equation}
where $\Cont(S,\Rdisc)$ is the \emph{locally constant} real functions. For $S$
the Cantor set, embedded in $\mathbb{R}$ in the standard way, the embedding is
continuous but not locally constant, so \eqref{eq:2:R-mod-Rdisc-general} is very
much nonzero. The cokernel is controlled by exactly these ``non-locally-constant''
real-valued functions on general light profinite sets.
\end{example}

Scholze closed by stating (with proofs deferred to the next lecture) the good
categorical properties that motivate working with \emph{light} condensed abelian
groups specifically.

\begin{theorem}[Properties of $\CondAbl$]
\label{thm:2:condab-properties}
The category $\CondAbl$ is a convenient homological setting. In particular:
\begin{enumerate}
\item It has all countable limits, and countable products are exact; a form of
Grothendieck's axiom \textup{AB6} holds for countable products (products
commuting suitably with filtered colimits). This is weaker than for all condensed
abelian groups, where \emph{all} products are exact; here only the countable ones
are.
\item The free light condensed abelian group on the convergent sequence,
$\mathbb{Z}[\mathbb{N}\cup\{\infty\}]$, is \emph{internally projective}.
\end{enumerate}
\end{theorem}

The forgetful functor $\CondAbl\to\CondSetl$ has a left adjoint $X\mapsto
\mathbb{Z}[X]$, the free light condensed abelian group; applied to a light
profinite set $S$ it gives the free object $\mathbb{Z}[S]$.

\begin{remark}[Internal projectivity is special to the light setting]
\label{rmk:2:internally-projective}
Scholze emphasized \Cref{thm:2:condab-properties}(2) as one of the main reasons
for the switch to the light setting. The category $\CondAbl$ no longer has enough
compact projective objects. Nevertheless $\mathbb{Z}[\mathbb{N}\cup\{\infty\}]$
--- which would \emph{not} be projective among all condensed abelian groups ---
happens to be projective, and even internally projective, in $\CondAbl$. This is
the only variant of condensed abelian groups in which Scholze is aware of any
nontrivial internally projective object: in all condensed abelian groups there
are plenty of projectives, but (apart from trivial cases such as $\mathbb{Z}$)
none are internally projective, and those that exist come from extremely
disconnected sets and are impractically large. Since the free group on a
convergent sequence is a very basic building block of the theory, it is
particularly valuable that in the light setting it enjoys these good categorical
properties. Scholze added that the failure of \emph{all} products to be exact was
one reason he was at first hesitant to make the switch.
\end{remark}
